
This study investigated whether CoT prompting enables models to leverage in-context uncertainty signals for confidence calibration. Systematic verification confirmed the mechanism chain: CoT prompting produces multi-step reasoning (100\% rate), reasoning contains epistemic hedging markers (82\% presence), markers structurally precede confidence generation (100\% compliance), and hedging count negatively correlates with verbalized confidence (r = -0.315, p < 1e-16).

These findings support a self-reading interpretation in which models incorporate uncertainty signals from their own reasoning into confidence judgments. The mechanism operates through prompting alone without model fine-tuning.

The primary limitation is that end-to-end ECE improvement was not measured; the ECE comparison experiment was not executed with real API calls. The evidence validates the mechanism linking CoT to hedging-confidence correlation but does not demonstrate that this correlation produces meaningful calibration improvement as measured by ECE.

\subsubsection{Future Work}

Four directions extend this work:

\begin{enumerate}
\item \textbf{Multi-model replication} on Llama-2-70B, Claude, and GPT-4 would establish generalization.
\item \textbf{Intervention study} manipulating hedging markers would provide causal evidence.
\item \textbf{Full ECE comparison} with real API calls would validate end-to-end calibration improvement.
\item \textbf{Token-padding control} would rule out length artifacts.
\end{enumerate}
